
\newpage
\section{Vector Fields and flows}

\vspace{0.4cm}\emph{\textbf{(Lecture 4)}}\\


Recall that a vector field $X\in \Gamma(TM)$ is a differential operator $X: C^\infty(M) \to C^\infty(M)$. In local coordinates $(x^1, \dots, x^n)$, a vector field can be expressed as $X = X^i \frac{\partial}{\partial x^i}$.

\begin{colordef}{Lie Bracket/Commutator of Vector Fields}{lie_bracket}
Let $X, Y \in \Gamma(TM)$. The Lie bracket (or commutator) $[-,-] : \Gamma(TM) \times \Gamma(TM) \to \Gamma(TM)$ of $X$ and $Y$ is defined by
\[
[X, Y] := X \circ Y - Y \circ X
\]
where $X$ and $Y$ are viewed as differential operators on $C^\infty(M)$.
\end{colordef}

Although $X \circ Y$ is a second order differential operator, $[X, Y]$ is a vector field (first order differential operator). Which can be seen by considering its coordinate expression.
\begin{align*}
    [X, Y]^h &= X(Y^h) - Y(X^h) \\
&= X^i \frac{\partial}{\partial x^i} (Y^h \frac{\partial}{\partial x^h}) - Y^i \frac{\partial}{\partial x^i} (X^h \frac{\partial}{\partial x^h}) \\
&= X^i \frac{\partial Y^h}{\partial x^i} - Y^i \frac{\partial X^h}{\partial x^i}.
\end{align*}

\begin{colremark}
Properties (see exercise 3.3):
\begin{itemize}
    \item $\mathbb{R}$ linearity
    \item $[X, Y] = -[Y, X]$ (anti-commutativity/-symmetry)
    \item Jacobi's identity
    \[
    [[X, Y], Z] + [[Z, X], Y] + [[Y, Z], X] = 0
    \]
    \item $[f X, g Y] = f \cdot g [X, Y] + X(g) \cdot f \cdot Y - Y(f) \cdot g \cdot X$
\end{itemize}
\end{colremark}

If $[X, Y]$ were a tensor field, then it would have to be of type (1,2). But $[fX,Y]\neq f[X,Y]$, as it involves derivatives of $f$. So $[-,-]$ is not tensorial.\\

% One can easily show that for $X,Y$ being coordinate vector fields, the Lie bracket vanishes, i.e. in any local coordinate charts
% \begin{equation}
%     \label{eq:coordinate_fields_commute}
%     \left[\frac{\partial}{\partial x^i}, \frac{\partial}{\partial x^j}\right] = 0.
% \end{equation}

% Indeed, let $(U,x)$ be a coordinate chart on a smooth manifold $M$. Recall that $\frac{\partial}{\partial x^i}(x^j) = \delta_i^j$ is constant. Hence
% \[
% \left[\frac{\partial}{\partial x^i}, \frac{\partial}{\partial x^j}\right](x^k)
% =
% \frac{\partial}{\partial x^i}\!\left(\frac{\partial x^k}{\partial x^j}\right)
% -
% \frac{\partial}{\partial x^j}\!\left(\frac{\partial x^k}{\partial x^i}\right)
% =0
% \]
% for all $k$, so coordinate vector fields commute.\\

One can easily show that for $X,Y$ being coordinate vector fields, the Lie bracket vanishes, i.e. in any local coordinate chart
\begin{equation}
    \label{eq:coordinate_fields_commute}
    \left[\frac{\partial}{\partial x^i}, \frac{\partial}{\partial x^j}\right] = 0.
\end{equation}

Indeed, let $(U,x)$ be a coordinate chart on a smooth manifold $M$. Recall that $\frac{\partial}{\partial x^i}(x^j) = \delta_i^j$ is constant. Hence, for any coordinate function $x^k$,
\[
\left[\frac{\partial}{\partial x^i}, \frac{\partial}{\partial x^j}\right](x^k)
=
\frac{\partial}{\partial x^i}\!\left(\frac{\partial}{\partial x^j}(x^k)\right)
-
\frac{\partial}{\partial x^j}\!\left(\frac{\partial}{\partial x^i}(x^k)\right)
=
\frac{\partial}{\partial x^i}(\delta_j^k)
-
\frac{\partial}{\partial x^j}(\delta_i^k)
=0.
\]

Thus the bracket vanishes on all coordinate functions and is therefore the zero field.\\

$(\Gamma(TM), [-,-])$ is an infinite-dimensional Lie algebra with the Lie group of diffeomorphisms of $M$ to itself, $\operatorname{Diff}(M)$, provided that $M$ is compact. (The Lie group is a manifold equipped with a compatible group structure such that the group operations are smooth maps.)



\subsection{Flow of a vector field}


Let $X \in \Gamma(TM)$.

\begin{colordef}{Integral Curve of a Vector Field}{}
A curve $\gamma: (a,b) \longrightarrow M$ is an integral curve of $X$ if
\[
\dot{\gamma}(t) = X\circ \gamma(t), \quad \forall t \in (a,b).
\]
\end{colordef}

In local coordinates:
\[
(\dot\gamma)^h(t) = X^h(\gamma(t)).
\]

We recall the following result from classic ODE theory.

\begin{colortheo}{Local Well-Posedness of ODEs}{}
Let $F:\mathbb{R}\times\mathbb{R}^n\to\mathbb{R}^n$ be smooth. For every $u_0\in\mathbb{R}^n$, the IVP
$$
u'(t)=F(t,u(t)), \quad u(0)=u_0
$$
admits a unique maximal solution $u:(T_-,T_+)\to\mathbb{R}^n$, where $T_\pm\in(0,\infty]$, such that:

\begin{itemize}
    \item Maximality: $u$ cannot be extended beyond $(T_-,T_+)$.
    \item Blow-up alternative: If $T_+<\infty$ (resp. $T_->-\infty$), then $|u(t)|\to\infty$ as $t\uparrow T_+$ (resp. $t\downarrow T_-$).
    \item Smooth dependence: For any compact $[a,b]\subset (T_-,T_+)$, the map $u_0 \mapsto u|_{[a,b]}$ is smooth.
\end{itemize}

\end{colortheo}

Characterisation of maximality: the blow-up alternative is equivalent to the statement that if
$$
\sup_{t \in [0,T_+)} |u(t)| < \infty,
$$
then $T_+ = +\infty$ (and similarly for $T_-$). In other words, a maximal solution can be extended past any finite time unless its norm becomes unbounded/as long as it remains bounded.


Moreover, $u(t) = u(t; u_0)$ depends smoothly on its arguments. This means that the \emph{flow map} $\Phi(t, u_0) := u(t; u_0)$ is smooth in both $t$ and $u_0$. However, the flow is not defined on a uniform time interval for all initial conditions $u_0$; rather, for each $u_0$, there exists some time interval $(-T_-(u_0), T_+(u_0))$ on which the flow is defined.\\

By working in local coordinates, the above theorem implies the following:

\begin{colortheo}{Flow of a Vector Field}{flow_vector_field}
Let $X \in \Gamma(TM)$. For every open, precompact subset $U \subset M$ (compact closure $\overline{U}$), there exist $\delta > 0$ and a unique smooth map
\[
\Phi: (-\delta, \delta) \times \overline{U} \longrightarrow M
\]
called the flow map, such that for all $q \in \overline{U}$, the curve $\gamma(t) = \Phi_t(q)$ is an integral curve of $X$ starting from $q$, i.e.,
\[
\left\{
\begin{array}{l}
\dfrac{\mathrm{d}}{\mathrm{d}t} \Phi_t(q) = X|_{\Phi_t(q)} \\[0.3cm]
\Phi_0(q) = q
\end{array}
\right.
\]

Moreover, if $X$ is compactly supported, then $U = M$ and $\delta = +\infty$.
\end{colortheo}

\begin{proof}
See \cite{tsakanikas2024dg2}, chapter 7.3.
\end{proof}




\begin{colremark}
Two integral curves cannot intersect (uniqueness of solutions of ODEs) $\implies \Phi_t$ is a diffeomorphism on its image.
\end{colremark}

Moreover, they form a semigroup, meaning that for all $t_1, t_2 \in (-\delta, \delta)$ and $q \in \overline{U}$, we have
\[
\Phi_{t_1+t_2}(q)
= (\Phi_{t_1} \circ \Phi_{t_2})(q).
\]


\begin{colexample}{}{}
    \label{ex:flow_of_coordinate_vector_field}
If $X = \frac{\partial}{\partial x^1}$, we solve the flow equation. The differential equation for the flow $\Phi_t(x) = (x^1(t), \ldots, x^n(t))$ is
\[
\frac{\mathrm{d}}{\mathrm{d}t} \Phi_t^k(x) = X^k(\Phi_t(x))
\]
where $X^k$ are the components of $X$. Since $X = \frac{\partial}{\partial x^1}$, we have $X^k = \delta_1^k$ and in particular $X^k(\Phi_t(x))=\delta^1_k$, so:
\[
\frac{\mathrm{d}}{\mathrm{d}t} \Phi_t^k(x) = \delta_1^k = \begin{cases} 1 & \text{if } k=1 \\ 0 & \text{if } k \neq 1 \end{cases}
\]

With initial condition $\Phi_0(x^1, \ldots, x^n) = (x^1, \ldots, x^n)$, integrating gives:
\[
\Phi_t((x^1, \ldots, x^n)) = (x^1 + t, x^2, \ldots, x^n).
\]

So the flow translates points in the $x^1$-direction at unit speed, as expected.
\end{colexample}



\begin{colorlem}{Straightening a Vector Field}{}
Let $X \in \Gamma(TM)$ with $X_p \neq 0$ at some point $p \in M$. Then there exists a local coordinate system around $p$ in which $X = \frac{\partial}{\partial x^1}$.
\end{colorlem}
\begin{proof}
See exercise 2.4.
\end{proof}


\begin{colorlem}{Commuting Vector Fields Generate Coordinates}{}
Let $X, Y \in \Gamma(TM)$ be smooth, $[X, Y] = 0$, and let $X_p$ and $Y_p$ be linearly independent at $p \in M$, then there exists local coordinates $(x^1, x^2)$ around $p$ such that $X = \frac{\partial}{\partial x^1}$ and $Y = \frac{\partial}{\partial x^2}$.
\end{colorlem}

\begin{proof}
See exercise 4.4.
\end{proof}



If we drop the compactness assumption, the flow from Theorem \ref{thm:flow_vector_field} doesn't necessarily exist for all times. To see this consider the following example.
\begin{colexample}{}{}

On $\mathbb{R}$: Let $X = y^2 \frac{\partial}{\partial y}$ and $y_0 > 0$. The integral curve $y(t)$ satisfies $\dot{y} = y^2$, which gives
\[
y(t) = \frac{1}{\frac{1}{y_0} - t}.
\]

This solution is only defined for $t < \frac{1}{y_0}$, showing that the flow reaches infinity in finite time.
\end{colexample}

A vector field generates the flow(s), a one-parameter family of diffeomorphisms. One can form the commutator on the group of such diffeomorphisms.\\

Recall now Definition \ref{def:lie_bracket}. 

\begin{colremark}
One can show that in local coordinates, the commutator of two flows can be expressed as
\[
[X,Y]^k = \lim_{t,s \to 0} \frac{\left(\Phi^{(Y)}_{-s} \circ \Phi^{(X)}_{-t} \circ \Phi^{(Y)}_{s} \circ \Phi^{(X)}_{t}(p)\right)^k - p^k}{ts}.
\]
\end{colremark}

This expression quantifies how much the flows of $X$ and $Y$ fail to commute. The four successive flows applied in sequence would return to $p$ if $[X,Y] = 0$; otherwise, we drift by an amount proportional to $ts$. The Lie bracket captures the leading-order term in this drift. Thus the geometric notion of flow non-commutativity coincides with the algebraic definition of the commutator.




\subsection{Lie derivative}
\begin{colordef}{Lie Derivative}{}
Let $X \in \Gamma(TM)$ and $\{\phi_t\}_{t \in \mathbb{R}}$ be its flow. The Lie derivative is defined by:
\begin{itemize}
    \item For $f \in C^\infty(M)$:
    \[
    \mathcal{L}_X f(p) = \lim_{t \to 0} \frac{f(\phi_t(p)) - f(p)}{t}
    \]
    \item For $Y \in \Gamma(TM)$:
    \[
    (\mathcal{L}_X Y)_p
    = \lim_{t \to 0} \frac{(\mathrm{d}\phi_{-t})_{\phi_t(p)}(Y|_{\phi_t(p)}) - Y|_p}{t}
    = \lim_{t \to 0} \frac{Y|_{\phi_t(p)} - \mathrm{d}(\phi_{+t})_p(Y|_p)}{t}
    \]
\end{itemize}
\end{colordef}

% \begin{figure}[H]
% \begin{center}
% \begin{tikzpicture}[scale=1]

% % Bottom curve
% \draw[thick] (0,0) .. controls (2,0) and (4,0) .. (6,0);

% % Middle curve
% \draw[thick] (0,1) .. controls (2,0.8) and (4,1.2) .. (6,1.8);

% % Top curve
% \draw[thick] (0,2) .. controls (2,1.8) and (4,2.2) .. (6,2.8);

% % X-field: bottom row of horizontal arrows
% \foreach \x in {0.8,2.2,3.6,5} {
%     \draw[thick,->] (\x,0.0) -- ++(0.8,0);
% }
% % Middle row of horizontal arrows tangent to the middle curve
% \foreach \x/\slope in {0.8/0.1,2.2/-0.1,3.6/0.15,5/0.2} {
%     \draw[thick,->] (\x,0.8) -- ++(0.7,\slope*0.7);
% }

% % Top row of horizontal arrows tangent to the top curve
% \foreach \x/\slope in {0.8/0.05,2.2/0.15,3.6/0.2,5/0.25} {
%     \draw[thick,->] (\x,1.8) -- ++(0.7,\slope*0.7);
% }

% % Y-field: row of vertical, upwards arrows
% \foreach \x/\h in {1.5/0.7,3/0.9,4.5/1.2,5.5/1.5} {
%     \draw[thick,->] (\x,0) -- (\x,\h);
% }

% % Label of Y-vector field
% \node at (5.8,0.9) {$Y$};

% % Label with left arrow
% \node at (9,0.9) {$\mathcal{L}_XY=0$};

% \end{tikzpicture}
% \end{center}
% \end{figure}


\begin{colortheo}{Lie Derivative of Functions and Vector Fields}{}
Let $X \in \Gamma(TM)$ with flow $\{\phi_t\}_{t \in \mathbb{R}}$. Then:
\begin{align*}
\mathcal{L}_X f &= X(f) \quad \text{for all } f \in C^\infty(M), \\
\mathcal{L}_X Y &= [X, Y] \quad \text{for all } Y \in \Gamma(TM).
\end{align*}
\end{colortheo}

\begin{proof}
    See handwritten notes from the lecture.
\end{proof}



\begin{colorlem}{Extension of Lie Derivative to Tensor Fields}{lie_derivative_extension_tensor}
The Lie derivative $\mathcal{L}_X$ extends to all tensor fields by requiring two \emph{axioms}:
\begin{enumerate}
    \item For $\omega \in \Gamma(T^*M)$: $(\mathcal{L}_X \omega)(Y) := \mathcal{L}_X(\omega(Y)) - \omega(\mathcal{L}_X Y)$
    \item Leibniz rule: $\mathcal{L}_X(T_1 \otimes T_2) = (\mathcal{L}_X T_1) \otimes T_2 + T_1 \otimes (\mathcal{L}_X T_2)$
\end{enumerate}

This yields a map
% \[
% \mathcal{L}_X: \Gamma(\otimes^k TM \otimes^l T^*M) \longrightarrow \Gamma(\otimes^k TM \otimes^l T^*M)
% \]
\[
\mathcal{L}: \Gamma(TM)\times T^k_l(M) \longrightarrow T^k_l(M)
\]
with $\mathcal{L}_X$ commuting with contractions.

\todo[inline]{maybe add the coordinate expression for this map}

\end{colorlem}

Note that this extension is unique, but we will not prove this here.

We remark that the first axiom is natural, because $\mathcal{L}_X(\omega(Y)) = (\mathcal{L}_X \omega)(Y) + \omega(\mathcal{L}_X Y)$ mimics a sort of chain rule.


\begin{proof}[Proof sketch]
Build inductively using the two axioms. For $(k,l)$-tensors, decompose as products of vectors and covectors, apply the Leibniz rule.
\end{proof}


Further more, this construction respects contractions (Definition \ref{def:contraction_tensors_coordinates}). Indeed, if $w \in \Gamma(\mathrm{T}^*M)$, $Y \in \Gamma(\mathrm{T}M)$, then $w \otimes Y$ is of type $(1,1)$, and in coordinates we have
$$
(w \otimes Y)^\alpha_\beta = w_\beta Y^\alpha.
$$

If one applies the contraction one gets
$$
\operatorname{tr}(w \otimes Y) = (w \otimes Y)^\alpha_\alpha = w_\alpha Y^\alpha.
$$

But
$$w(Y) = w_\beta dx^\beta \left(Y^\alpha \frac{\partial}{\partial x^\alpha}\right)=w_\alpha Y^\alpha.$$

And therefore
\begin{equation}
    \label{eq:lie_deriv_contract1}
    \operatorname{tr}(w \otimes Y) =w(Y).
\end{equation}

Now apply $\mathcal{L}_X$ to both sides and apply the first axiom/requirement,
\begin{align*}
    \mathcal{L}_X(\operatorname{tr}(w \otimes Y))
    &= \mathcal{L}_X(w(Y))\\
    &\stackrel{1.}{=} (\mathcal{L}_X w)(Y) + w(\mathcal{L}_X Y)
\end{align*}

Now we use the second axiom/requirement and apply the contraction on both sides.
\begin{align*}
\operatorname{tr}(\mathcal{L}_X(w \otimes Y))
&\stackrel{2.}{=} \operatorname{tr}((\mathcal{L}_X w) \otimes Y) + \operatorname{tr}(w \otimes (\mathcal{L}_X Y)) \\
&\stackrel{\eqref{eq:lie_deriv_contract1}}{=} (\mathcal{L}_X w)(Y) + w(\mathcal{L}_X Y)
\end{align*}

That way, comparing the two expressions, we have established the following lemma.

\begin{colorlem}{Lie Derivative Commutes with Contraction for $(1,1)$ Tensors}{}
Let $w \in \Gamma(\mathrm{T}^*M)$ and $Y \in \Gamma(\mathrm{T}M)$. Then the Lie derivative commutes with contraction:
\[
\mathcal{L}_X(\operatorname{tr}(w \otimes Y)) = \operatorname{tr}(\mathcal{L}_X(w \otimes Y)).
\]
\end{colorlem}


One can also construct this extension the other way round, i.e. we require the Leibnitz rule + the contraction compatibility, and we get the first axiom as a consequence.
\begin{align*}
\mathcal{L}_X w(Y) 
&\stackrel{\eqref{eq:lie_deriv_contract1}}{=} \mathcal{L}_X(\operatorname{tr}(w \otimes Y)) \\
&= \operatorname{tr}(\mathcal{L}_X(w \otimes Y)) \\
&= \operatorname{tr}(\mathcal{L}_X w \otimes Y) + \operatorname{tr}(w \otimes \mathcal{L}_X Y)\\
&= (\mathcal{L}_X w)(Y) + w(\mathcal{L}_X Y)
\end{align*}



\begin{colremark}
\label{rem:lie_derivative_extension_tensor_fields}
The results from before can be summarised as follows. We have 4 statements.
\begin{enumerate}[itemsep=0pt, topsep=0pt]
    \item $(\mathcal{L}_X \omega)(Y) := \mathcal{L}_X(\omega(Y)) - \omega(\mathcal{L}_X Y)$
    \item Leibnitz rule for tensor products
    \item $\mathcal{L}_X$ commutes with contractions
    \item Extension of $\mathcal{L}_X$ to all tensor fields\\
\end{enumerate}

Then we have
$$
1.\text{ and } 2.
\iff
3.\text{ and } 4.
\iff
(2.\text{ and } 3. \implies 1.)
$$
\end{colremark}


\begin{colremark}
In general: Tensors evaluated at their arguments, like $w(Y)$, $g(X,Y)$, can be viewed as contractions of tensor products. E.g.:

$$
g(X,Y) = g_{ab} X^a Y^b = \operatorname{tr}(\operatorname{tr}(g \otimes X \otimes Y))
$$
\end{colremark}


\begin{colremark}
\label{rem:lie_derivative_tensor_coordinates}
If $T \in T^k_l(M)$, and $X = \frac{\partial}{\partial x^i}$, in local coordinates then $$(\mathcal{L}_X T)^{i_1 \ldots i_k}_{j_1 \ldots j_l} = \frac{\partial T^{i_1 \ldots i_k}_{j_1 \ldots j_l}}{\partial x^i}.$$
See the solution of exercise 4.3b.
\end{colremark}



\begin{colorlem}{Lie Derivative Commutator}{lie_derivative_commutator}
For any $X, Y \in \Gamma(TM)$ and tensor field $T \in T^k_l(M)$, the Lie derivatives with respect to $X$ and $Y$ satisfy
\[
\mathcal{L}_{[X,Y]} T=\mathcal{L}_X \mathcal{L}_Y T - \mathcal{L}_Y \mathcal{L}_X T.
\]
\end{colorlem}

\begin{proof}
See exercise 4.3a.
\end{proof}




\begin{colordef}{Killing Vector Field}{killing}
Let $(M,g)$ be a Riemannian manifold. A vector field $X \in \Gamma(\mathrm{T}M)$ is called a \emph{Killing vector field} if
\[
\mathcal{L}_X g = 0.
\]
\end{colordef}

Remark: this is named after Wilhelm Killing, and not after literally \emph{killing} a vector field.\\


% For example, as seen in exercise 4.3b, if $X = \frac{\partial}{\partial x^1}$, then $X$ is a Killing vector field for $g$ if and only if $\frac{\partial g_{ij}}{\partial x^1} = 0$, i.e., the metric components are independent of $x^1$. In this case, the flow of $X$ is an isometry. 
% \todo[inline]{ this is not really what 4.3b shows }

\begin{colorlem}{Isometric flows and Killing Vector Fields}{killing_generates_isometries}
Let $X \in \Gamma(\mathrm{T}M)$ with flow $\{\Phi_t\}$. If $X$ generates a flow of isometries, it is Killing, i.e.
$$(\Phi_t)^* (g|_{\Phi_t}) = g\implies \mathcal{L}_X g = 0.$$
\end{colorlem}
\begin{proof}
See exercise 4.3b.
\end{proof}

We will not show the converse in full generality, but we proved a related statement in exercise 5.1. There we show that if a metric is independent of a coordinate $x^k$, then the flow $F_t$ of $\frac{\partial}{\partial x^k}$ is an isometry, $F^*_t(g|_{F_t}) = g$. We use this several times, for example in exercise 7.2b.





\subsection{Connections}

Lie derivatives measure how a tensor field changes along the flow of a vector field. For vector fields $V$ and $Y$, $\mathcal{L}_V Y$ is not suitable as a directional derivative because it is not $\mathbb{R}$-linear in the first argument: for $f \in C^\infty(M)$,
\[
\mathcal{L}_{fV} Y \neq f\,\mathcal{L}_V Y ,
\]
so $\mathcal{L}_V Y|_p$ depends on $V$ near $p$, not just on $V|_p$.

Given $Y \in \Gamma(TM)$, $p \in M$, and $V \in T_pM$, we want a directional derivative $D_V Y|_p$ that depends only on $V|_p$. This requires comparing tangent vectors at different points, which is not possible immediately, as they live in different spaces. A systematic way to do this defines a \emph{connection}, enabling differentiation of vector fields, curve acceleration, and parallel transport.




\begin{colordef}{Connection}{connection}
A connection $\nabla$ on $M$ is a map
$$\nabla \colon \Gamma(\mathrm{T}M) \times \Gamma(\mathrm{T}M) \to \Gamma(\mathrm{T}M)$$
such that:
\begin{enumerate}
    \item $\nabla$ is $C^\infty(M)$-linear in the first argument:
    \[
    \nabla_{X_1 + f \cdot X_2} Y = \nabla_{X_1} Y + f \cdot \nabla_{X_2} Y
    \]
    \item $\nabla$ is $\mathbb{R}$-linear in the second argument:
    \[
    \nabla_X (Y_1 + a Y_2) = \nabla_X Y_1 + a \nabla_X Y_2, \quad a \in \mathbb{R}
    \]
    \item Satisfies the Leibniz rule with respect to the second argument:
    \[
    \nabla_X (f Y) = X(f) \cdot Y + f \cdot \nabla_X Y
    \]
\end{enumerate}
It is also called \emph{covariant derivative}.
\end{colordef}

$\nabla Y$ is a $(1,1)$-tensor field. It is also local: if $U \subseteq M$ is open and $p \in U$, then $(\nabla_X Y)|_p$ depends only on $X|_p$ and $Y|_U$.\\

For the existence of connections, see Theorem \ref{thm:levi_civita}.

Because of the Leibnitz rule in the second argument, connections are not $C^\infty(M)$-linear in the second argument (there is this extra term $X(f)\cdot Y$) and therefore they are not tensor fields.

E.g. on $\mathbb{R}^n$ one can verify that the \emph{Euclidean connection} defined via the standard directional derivative $D$
\begin{equation}
    \label{eq:connection_Rn}
    (D_X Y)^k = X^i \frac{\partial Y^k}{\partial x^i}.
\end{equation}
is in fact a connection.

\begin{colorlem}{}{}
If $\nabla$ is a connection on $M$, it defines a unique connection on any open subset $U \subseteq M$.
\end{colorlem}


\begin{proof}
    Let $X, Y \in \Gamma(\mathrm{T}U)$, $p \in U$. We want to show that $\nabla_X Y|_p$ is uniquely defined.

    Let $F: M \to [0,1]$ be a smooth cutoff function such that
    $$
    F(q)=\begin{cases}
    1 \text{ if $q$ is in a neighborhood of } p \text{ in } U \\
    0 \text{ on } M \setminus U.
    \end{cases}
    $$

    Then $\widetilde{X} = F \cdot X$, $\widetilde{Y} = F \cdot Y$ extend to the whole of $M$.

    Define $\nabla_X Y|_p = \nabla_{\widetilde{X}} \widetilde{Y}|_p$. This is independent of the choice of $F$: If $Y$, $\widetilde{Y}$ are two vector fields such that $\widetilde{Y} = Y$ in a neighbourhood $V$ of $p$ and if $f: M \to [0,1]$,
    \[
    f =
    \begin{cases}
        1 & \text{on } V' \subset V \\
        0 & \text{on } M \setminus V
    \end{cases}
    \]
    then
    $$
    \left. \nabla_X \left( f \cdot (Y - \widetilde{Y}) \right) \right|_p = \left. X(f) \right|_p \cdot (\underbrace{Y - \widetilde{Y}}_{=0}) \big|_p + \underbrace{\left. f \right|_p}_{=1} \cdot (\left. \nabla_X (Y - \widetilde{Y}) \right|_p)
    $$

    But $f \cdot (Y - \widetilde{Y}) = 0$ in a neighborhood of $p$ $\implies$ $\left. \nabla_X Y \right|_p = \left. \nabla_X \widetilde{Y} \right|_p$

\end{proof}

We will use the above in order to talk about $\nabla_X Y$ when $X, Y$ are not defined on the whole of $M$ (e.g. coordinate vector fields), or fields on a chart.

\begin{colordef}{Christoffel Symbols}{christoffel}
Let $(x^1, \ldots, x^n)$ be a local coordinate system. The Christoffel symbols $\Gamma_{ij}^k$ are defined by
\[
\nabla_{\frac{\partial}{\partial x^i}} \frac{\partial}{\partial x^j} = \Gamma_{ij}^k \frac{\partial}{\partial x^k}.
\]
\end{colordef}
Equivalently, 
$$\Gamma_{ij}^k = \mathrm{d}x^k\!\left( \nabla_{\frac{\partial}{\partial x^i}} \frac{\partial}{\partial x^j} \right) = \left(\nabla_{\frac{\partial}{\partial x^i}} \frac{\partial}{\partial x^j}\right)^k$$ 
is the $k^{\text{th}}$ component.

Note: $\Gamma_{ij}^k$ are not components of a tensor (see exercise 4.1a). If they would, having zero Christoffel symbols in one coordinate system (say cartesian) would imply they vanish in all coordinate systems, which is not the case.\\

In local coordinates, the covariant derivative of a vector field $Y$ in the direction of $X$ is given by
\begin{equation}
    \label{eq:covariant_derivative_coordinates}
    (\nabla_X Y)^k = X^i \frac{\partial Y^k}{\partial x^i} + \Gamma_{ij}^k X^i Y^j.
\end{equation}


\todo[inline]{ add a clear example comparing this to the lie derivative }



\begin{colorlem}{Extension of Connections to Tensor Fields}{}
Let $M$ be a smooth manifold equipped with a connection $\nabla$. The connection extends to a map
\[
\nabla: \Gamma(TM) \times T^k_l(M) \longrightarrow T^k_l(M)
\]
by requiring:
\begin{enumerate}
    \item The Leibniz rule for tensor products: $\nabla_X(S \otimes T) = (\nabla_X S) \otimes T + S \otimes (\nabla_X T)$
    \item Compatibility with contractions: $\nabla_X(\operatorname{tr} A) = \operatorname{tr}(\nabla_X A)$
\end{enumerate}
\end{colorlem}

The coordinate formula for $X=\frac{\partial}{\partial x^a}$ is
\[
(\nabla_{\partial_a} T)^{i_1 \dots i_k}_{j_1 \dots j_l}
= \partial_a T^{i_1 \dots i_k}_{j_1 \dots j_l}
+ \sum_{r=1}^{k} \Gamma^{i_r}_{a b} \, T^{i_1 \dots b \dots i_k}_{j_1 \dots j_l}
- \sum_{s=1}^{l} \Gamma^{b}_{a j_s} \, T^{i_1 \dots i_k}_{j_1 \dots b \dots j_l}
\]


\todo[inline]{ $\nabla_X$ for general $X$? (the example below states it for 1,1 tensors)}
Compare to extension of the Lie derivative Lemma \ref{lem:lie_derivative_extension_tensor} and Remark \ref{rem:lie_derivative_extension_tensor_fields}.


Compare to the coordinate expression of the Lie derivative in Remark \ref{rem:lie_derivative_tensor_coordinates}.


\todo[inline]{ explain how to compute Lie and covariant derivatives of tensors, not using the large formula, but using the axioms/requirements, which is more intuitive and easier }


\begin{proof}

Recall equation \eqref{eq:covariant_derivative_coordinates}. To derive a formula for general $k,l$ tensors, we first extend $\nabla$ to 1-forms. One can use the contraction axiom to get
$$(\nabla_X \omega)(Y) = X(\omega(Y)) - \omega(\nabla_X Y)$$

which yields the coordinate formula
\[
(\nabla_{\partial_i} \omega)_j = \partial_i \omega_j - \Gamma^k_{ij} \omega_k
\]

Because $\nabla_X$ is $C^\infty(M)$-linear in the vector field argument and linear over sums in tensors, knowing how it acts on the coordinate basis fully determines its action on any tensor field.

See Exercise 4.2 for the full proof.
\end{proof}


\begin{colexample}{}{}

For a (0,1)-tensor, i.e. a 1-form:
\[
(\nabla_X \omega)_k = X^i \frac{\partial \omega_k}{\partial x^i} - \Gamma_{ik}^{\alpha} X^i \omega_\alpha
\]

For $(1,1)$ tensor fields:

\[
(\nabla_X T)^a_b = X^i \frac{\partial T^a_b}{\partial x^i} + \Gamma_{i\alpha}^{a} X^i T^\alpha_b - \Gamma_{ib}^{m} X^i T^a_m
\]

\end{colexample}




\begin{colordef}{Torsion of a Connection}{}
Let $\nabla$ be a connection on $M$. The torsion tensor $T$ is defined by
\[
T(X, Y) = \nabla_X Y - \nabla_Y X - [X, Y]
\]
for all $X, Y \in \Gamma(TM)$. Here the term $\nabla_Y X + [X, Y]$ is called \emph{adjoint torsion}.
\end{colordef}

$T$ is an anti-symmetric $(1,2)$ tensor, see exercise 4.1b.

\begin{colremark}[Torsion for coordinate vector fields]

For coordinate vector-fields,
$$\left[\frac{\partial}{\partial x^i}, \frac{\partial}{\partial x^j}\right] = 0$$


so in view of Definition \ref{def:christoffel}, we have
$$T_{ij}^k = \Gamma_{ij}^k - \Gamma_{ji}^k,$$

and there fore also
$$
T = 0 \iff \Gamma_{ij}^k = \Gamma_{ji}^k.
$$

Suffices to be true in one coordinate system.

\end{colremark}
